Die Photosynthese ist der wichtigste Prozess auf der Erde, der durch Nahrung, Sauerstoff und andere Ressourcen Leben erhält. Diese Bedeutung bringt jedoch auch Herausforderungen mit sich. Es handelt sich um einen komplexen Prozess mit vielen interagierenden Komponenten und Rückkopplungsschleifen. Diese Komplexität hat die Menschen jedoch nicht davon abgehalten, zu versuchen, ihn zu verstehen. In den letzten Jahrhunderten stand die Photosynthese im Mittelpunkt der biologischen Forschung. Seitdem hat sich der Umfang der wissenschaftlichen Arbeit jedoch stabilisiert. Um die Möglichkeiten weiter zu erforschen, müssen mehr Feinheiten und Details des Prozesses verstanden werden, und eines der gängigsten Werkzeuge dafür ist die mathematische Modellierung. Diese Modelle sind keine neuen Erfindungen, aber sie haben im Laufe der Jahre an Popularität und Komplexität gewonnen. Der aktuelle Stand der Technik ist jedoch nicht ohne Mängel. Aufgrund der Komplexität der Photosynthese schränken Modelle ihren Umfang ein, um sie besser handhabbar zu machen. Diese Reduzierung führt jedoch zu vielen möglichen Modellen. Im Laufe der Jahre hat die Zahl der Veröffentlichungen, in denen Photosynthesemodelle erwähnt werden, stetig zugenommen. Diese Zunahme der Modelle bringt jedoch auch eine Komplikation mit sich: die Gefahr, von der Auswahl überwältigt zu werden. Es gab Versuche, dieses Problem zu lösen, aber sie waren nur teilweise erfolgreich. Diese Lösungen zielen darauf ab, Probleme über möglichst viele Modelle hinweg zu lösen, aber genau das untergräbt ihre Leistungsfähigkeit. Die Modelle unterscheiden sich so stark voneinander, dass viele Abkürzungen erforderlich sind, um sie alle in einen einzigen Rahmen zu integrieren. Umgekehrt verringert die Berücksichtigung der Modellvielfalt und die Konzentration auf bestimmte Modelltypen den Umfang der Lösung. Auf diese Weise kann das Augenmerk auf bestimmte Modelltypen gelenkt werden. Es kann mehr Sorgfalt darauf verwendet werden, sicherzustellen, dass diese Modelle einem Standard entsprechen, der modernen Richtlinien wie \glsentryshort{fair} und \glsentryshort{cure} folgt. Darüber hinaus ermöglicht dies einen direkteren Vergleich, sodass standardisierte Demonstrationen verwendet werden können, um die Fähigkeiten der Modelle zu präsentieren. Darüber hinaus trägt die Bereitstellung einer Webplattform zum Austausch dieser Modelle mit der Community dazu bei, dass sie leichter gefunden und genutzt werden können. Diese drei Aspekte bilden den Kern von \greensloth{}, das ein ganzes Ökosystem von Tools zum Erstellen, Demonstrieren und Teilen von \glsentryshort{ode}-Modellen der Photosynthese umfasst. Die Website (\url{https://greensloth.rwth-aachen.de/}) ist die Hauptplattform für den Austausch der Modelle, und die Tools zu ihrer Erstellung und Demonstration sind auch als Open-Source-GitHub-Repositorys (\url{https://github.com/ElouenCorvest/GreenSloth/tree/main}) verfügbar. Um die Fähigkeiten von \greensloth{} hervorzuheben, werden fünf Modelle der Photosynthese verwendet, um die Stärken und Schwächen des Ökosystems aufzuzeigen. Diese Modelle unterscheiden sich hinsichtlich des Umfangs der Photosynthese und ihrer Komplexität und bieten einen guten Ausgangspunkt als Proof of Concept.